\section{Method}\label{sec:method}

This section describes CIDA as the paper and its supplementary material specify it~\citep{kundu2020cida}.
\Cref{sec:implementation} then lists the details we had to choose ourselves and the changes we made.

% Figure 1: CIDA architecture (simplified from docs/cida_architecture.drawio).
\begin{figure*}[t]
\centering
\begin{tikzpicture}[
  font=\small,
  >={Stealth[length=2mm]},
  net/.style={draw, rounded corners=2pt, minimum width=9mm, minimum height=6mm, fill=blue!10},
  frz/.style={draw, dashed, rounded corners=2pt, minimum width=9mm, minimum height=6mm, fill=gray!20},
  sp/.style={draw, ellipse, minimum width=15mm, minimum height=10mm, fill=orange!12},
  io/.style={inner sep=1pt, font=\footnotesize, align=center},
  loss/.style={font=\footnotesize, text=red!70!black},
  note/.style={font=\scriptsize, text=black!70},
  ttl/.style={font=\small\bfseries, anchor=west},
  ar/.style={->, thick},
]
% ---------------------------------------------------------------- (a) Stage 1
\node[ttl] at (-0.3,2.05) {(a) Stage 1: source training};
\node[io] (xs) at (0,0) {$\x_s$};
\node[net] (fs) at (1.1,0) {$f_s$};
\node[sp] (U) at (2.75,0) {$\U$};
\node[note] at (2.75,-0.72) {prototypes $\Ps^c$, $\Ps$};
\node[net] (gs) at (4.4,0) {$g_s$};
\node[io] (os) at (5.55,0) {$\Cs\cup\{\unk\}$};
\node[io] (neg) at (2.75,-1.35) {negatives $\uu_n\sim\Ps$, label $\unk$};
\draw[ar] (xs) -- (fs);
\draw[ar] (fs) -- (U);
\draw[ar] (U) -- (gs);
\draw[ar] (gs) -- (os);
\draw[ar] (neg.east) -| (gs.south);
\node[loss, above=1mm of U] {$\mathcal{L}_{s1}$};
\node[loss, above=1mm of gs] {$\mathcal{L}_{s2}$};

% divider
\draw[black!30, dashed] (6.45,-1.6) -- (6.45,2.2);

% ---------------------------------------------------------------- (b) Stage 2
\begin{scope}[xshift=7.6cm]
\node[ttl] at (-0.9,2.05) {(b) Stage 2: source-free adaptation};
\node[io] (px) at (0,0.95) {proxy-source\\$\uu_s^c\sim\Ps^c$};
\node[net] (fe) at (1.5,0.95) {$f_e$};
\node[io] (xt) at (0,-0.45) {$\x_t$, one-shot $\tilde{\x}_t^c$};
\node[net] (ft) at (1.5,-0.95) {$f_t$};
\node[note] at (1.5,-1.45) {initialised from $f_s$};
\node[sp, minimum width=17mm, minimum height=13mm, align=center] (V) at (3.3,0.0) {$\V$\\[-1pt]{\scriptsize guides $\vg^c$}};
\node[net] (fd) at (5.0,0.95) {$f_d$};
\node[frz] (gs2) at (6.2,0.95) {$g_s$};
\node[note] at (6.2,0.43) {frozen};
\node[net] (gt) at (5.6,-0.95) {$g_t$};
\node[io] (yh) at (8.3,0.0) {$\hat{y}$\\{\scriptsize $|\Ct|$ logits}};
\draw[ar] (px) -- (fe);
\draw[ar] (fe.east) -- (V.north west);
\draw[ar] (xt.south) |- (ft.west);
\draw[ar] (ft.east) -- (V.south west);
\draw[ar] (V.north east) -- (fd.west);
\draw[ar] (fd) -- (gs2);
\draw[ar] (V.south east) -- (gt.west);
\draw[ar] (gs2.east) -- node[above, note, pos=0.55] {$\Cs$} (yh.north west);
\draw[ar] (gt.east) -- node[below, note, pos=0.55] {$\Cp$} (yh.south west);
\node[loss, above=1mm of fd] {$\mathcal{L}_{r2}$: $f_d(f_e(\uu))\approx\uu$};
\node[loss, above=0.5mm of V] {$\mathcal{L}_{a1}$, $\mathcal{L}_{a2}$};
\node[loss, below=0.5mm of yh] {$\mathcal{L}_{r1}$, $\mathcal{L}_{c}$};
\end{scope}
\end{tikzpicture}
\caption{CIDA architecture. The source data never enter Stage~2, where Gaussian prototypes replace them.
(a)~Stage~1 trains $f_s$ and $g_s$. $\mathcal{L}_{s1}$ makes each class a compact cluster around its prototype, and $\mathcal{L}_{s2}$ teaches $g_s$ to assign synthetic negatives to the extra class $\unk$.
(b)~Stage~2 trains a target extractor $f_t$, a classifier $g_t$ for the new classes and two projections $f_e$ and $f_d$ between $\U$ and $\V$, while $g_s$ stays frozen.
Red labels show where each loss is computed, and \cref{tab:losses} lists the networks each loss updates.}
\label{fig:arch}
\end{figure*}


\subsection{Problem setting}\label{sec:setting}
The labelled source set is $\Ds=\{(\x_s,y_s)\}$ with label set $\Cs$.
The unlabelled target set $\Dt=\{\x_t\}$ is drawn from the marginal distribution $q_X$ and has a larger label set $\Ct\supset\Cs$.
The target-private classes are $\Cp=\Ct\setminus\Cs$, and each class $c\in\Cp$ comes with one labelled target image $\tilde{\x}_t^c$.
The source data are not available during adaptation.

\subsection{Stage 1: source training}\label{sec:stage1}
A feature extractor $f_s$ maps images to a feature space $\U$, and a classifier $g_s$ maps these features to $|\Cs|+1$ outputs.
The extra output, denoted $\unk$, stands for ``none of the source classes''.

\paragraph{Gaussian prototypes.}
For each class $c\in\Cs$, a Gaussian $\Ps^c=\mathcal{N}(\bm{\mu}_s^c,\bm{\Sigma}_s^c)$ is fitted to the features $f_s(\x_s)$ of that class.
A global Gaussian $\Ps=\mathcal{N}(\bm{\mu}_s,\bm{\Sigma}_s)$ is fitted to all source features in the same way.
These prototypes are all that Stage~2 keeps of the source data.

\paragraph{Class separability ($\mathcal{L}_{s1}$).}
This loss is a cross-entropy in which the class prototypes act as the classifier.
With $\uu_s=f_s(\x_s)$, it reads
\begin{equation}
  \mathcal{L}_{s1} = \E_{(\x_s,y_s)}\Big[-\log\frac{\exp \ell^{y_s}(\uu_s)}{\sum_{c\in\Cs}\exp \ell^{c}(\uu_s)}\Big],
  \label{eq:ls1}
\end{equation}
where $\ell^c(\uu)$ scores $\uu$ under $\Ps^c$.
The paper uses the density $\Ps^c(\uu)$ as this score, while we use the log-density $\log\Ps^c(\uu)$ for numerical reasons (\cref{sec:choices}).
Minimising the loss makes each class a compact cluster around its prototype.

\paragraph{Negative samples ($\mathcal{L}_{s2}$).}
Negative features $\uu_n$ represent regions of $\U$ where the source density is low.
To create one, we draw a candidate $\uu\sim\Ps$ from the global Gaussian and find its most likely class $\hat c=\arg\max_c \Ps^c(\uu)$.
Let $\lambda$ be the largest eigenvalue of $\bm{\Sigma}_s^{\hat c}$ and $\bm{\nu}$ its eigenvector.
The candidate is kept as a negative if
\begin{equation}
  \Ps^{\hat c}(\uu) < \Ps^{\hat c}\big(\bm{\mu}_s^{\hat c}+3\sqrt{\lambda}\,\bm{\nu}\big),
  \label{eq:neg}
\end{equation}
that is, if it lies outside the $3\sigma$ region of class $\hat c$.
The classifier then learns to assign these negatives to $\unk$ through
\begin{equation}
  \mathcal{L}_{s2} = \E_{(\x_s,y_s)}\,\lce\big(g_s(f_s(\x_s)),y_s\big) + \E_{\uu_n}\,\lce\big(g_s(\uu_n),\unk\big),
  \label{eq:ls2}
\end{equation}
where each batch holds as many negatives as source images.

\paragraph{Optimisation.}
One warm-up epoch first trains $f_s$ and $g_s$ with plain cross-entropy, which gives the first prototypes.
Training then alternates between the two losses, one per iteration.
$\mathcal{L}_{s1}$ updates only $f_s$, while $\mathcal{L}_{s2}$ updates both $f_s$ and $g_s$.
Each loss has its own Adam optimiser~\citep{kingma2015adam} with learning rate $10^{-4}$, and the prototypes are refitted after every epoch.
The final source model consists of $f_s$, $g_s$ and the prototypes $\{\Ps^c\}$ and $\Ps$.

\subsection{Stage 2: adaptation to the target}\label{sec:stage2}

\paragraph{Networks.}
The target feature extractor $f_t$ is initialised from $f_s$ and maps images to a new space $\V$.
A new classifier $g_t$ has one output per class in $\Cp$, while the source classifier $g_s$ stays frozen.
Two domain projections, $f_e:\U\to\V$ and $f_d:\V\to\U$, connect the two feature spaces.
Stage~2 replaces the source data with \emph{proxy-source} samples, which are features drawn from the prototypes, $\uu\sim\Ps^c$.
Before adaptation, $f_e$ and $f_d$ are pretrained towards the identity by minimising $\|f_e(\uu)-\uu\|^2$ and $\|f_d(\uu)-\uu\|^2$ on these samples.
For a feature $\fv\in\V$, the model then outputs $|\Ct|$ logits,
\begin{equation}
  \hat{y}(\fv) = \big[\,g_s(f_d(\fv))\big|_{\Cs}\;;\; g_t(\fv)\,\big],
  \label{eq:yhat}
\end{equation}
where the first part keeps the $|\Cs|$ shared-class logits of $g_s$ and drops $\unk$.

\paragraph{Retaining the source classes ($\mathcal{L}_{r1}$, $\mathcal{L}_{r2}$).}
Proxy-source samples $\uu_s^c\sim\Ps^c$ with $c\in\Cs$ enter two losses that preserve the source classes,
\begin{align}
  \mathcal{L}_{r1} &= \E_{c,\,\uu_s^c}\,\lce\big(\hat{y}(f_e(\uu_s^c)),c\big),\label{eq:lr1}\\
  \mathcal{L}_{r2} &= \E_{c,\,\uu_s^c}\,\big\|f_d(f_e(\uu_s^c))-\uu_s^c\big\|^2.\label{eq:lr2}
\end{align}
$\mathcal{L}_{r1}$ keeps the source classes recognisable after the round trip through $f_e$ and $f_d$.
$\mathcal{L}_{r2}$ keeps $f_d$ an inverse of $f_e$, which stops the projections from discarding information.

\paragraph{Guides and pseudo-labels.}
Each class $c\in\Ct$ has a guide $\vg^c\in\V$, which acts as its class centre in the target space,
\begin{equation}
  \vg^c = f_e(\bm{\mu}_s^c)\ \ \text{for } c\in\Cs,
  \qquad
  \vg^c = f_t(\tilde{\x}_t^c)\ \ \text{for } c\in\Cp.
  \label{eq:guides}
\end{equation}
The shared guides can move during adaptation, because $f_e$ is trained.
Each target sample receives the label of its nearest guide, $k(\x_t)=\arg\min_{c}\|f_t(\x_t)-\vg^c\|_2$.
Gaussian prototypes $\Pt^c$ are then fitted to the pseudo-labelled target features.
The confident set $\Bt^c$ holds the $n=20\%$ of samples labelled $c$ that lie closest to $\vg^c$.
The pseudo-labels, $\Pt^c$ and $\Bt^c$ are all recomputed after every epoch.

\paragraph{Alignment ($\mathcal{L}_{a1}$, $\mathcal{L}_{a2}$) and classification ($\mathcal{L}_{c}$).}
With $\fv_t=f_t(\x_t)$, the remaining three losses are
\begin{align}
  \mathcal{L}_{a1} &= \E_{\x_t\sim q_X}\Big[-\log\frac{\exp\ell_t^{k(\x_t)}(\fv_t)}{\sum_{c\in\Ct}\exp\ell_t^{c}(\fv_t)}\Big],\label{eq:la1}\\
  \mathcal{L}_{a2} &= \E_{c,\,\x_t\in\Bt^c}\,\|\fv_t-\vg^c\|^2,\label{eq:la2}\\
  \mathcal{L}_{c}  &= \E_{c,\,\x_t\in\Bt^c}\,\lce\big(\hat{y}(\fv_t),c\big),\label{eq:lc}
\end{align}
where $\ell_t^c$ is the score under $\Pt^c$.
$\mathcal{L}_{a1}$ separates the target clusters and makes room for the new classes.
$\mathcal{L}_{a2}$ pulls confident samples towards their guides, and $\mathcal{L}_{c}$ trains the classifiers on the confident pseudo-labels.

\paragraph{Optimisation.}
Each loss has its own Adam optimiser and updates only some of the networks, as \cref{tab:losses} shows.
The losses take turns in a fixed cycle, with one loss per iteration.
Because separate optimisers scale the gradients of each loss independently, the paper needs no loss weights.

\begin{table}[t]
  \centering
  \small
  \caption{Stage~2 losses and the networks each one updates. The source classifier $g_s$ stays frozen throughout.}
  \label{tab:losses}
  \begin{tabular}{@{}lll@{}}
    \toprule
    Loss & Samples & Updates \\
    \midrule
    $\mathcal{L}_{r1}$ \eqref{eq:lr1} & proxy-source $\uu_s^c\sim\Ps^c$ & $f_e, f_d, g_t$ \\
    $\mathcal{L}_{r2}$ \eqref{eq:lr2} & proxy-source $\uu_s^c\sim\Ps^c$ & $f_e, f_d$ \\
    $\mathcal{L}_{c}$ \eqref{eq:lc}   & confident set $\Bt$              & $f_t, g_t$ \\
    $\mathcal{L}_{a1}$ \eqref{eq:la1} & all target samples               & $f_t$ \\
    $\mathcal{L}_{a2}$ \eqref{eq:la2} & confident set $\Bt$              & $f_e, f_t$ \\
    \bottomrule
  \end{tabular}
\end{table}

\paragraph{Prediction.}
The predicted class of a target image $\x$ is $\arg\max_{c\in\Ct}\hat{y}(f_t(\x))_c$.
In the paper's few-shot setting, with several labelled images per private class, each private guide is the mean feature of those images.
